\subsection{発展的な読み物}
SWEBOK 02 章で紹介される発展的な読み物を、日本語で読む目的に合わせて整理する。

\par\smallskip\noindent\refstepcounter{table}\label{tab:ch02-08}表 \thetable: 発展的な読み物における文献・読むと得られること\par\smallskip
表 \thetable は、発展的な読み物における文献・読むと得られることを比較・確認しやすい形で整理したものである。\par\smallskip
\begin{longtable}{P{0.32\linewidth}P{0.60\linewidth}}
\toprule
文献 & 読むと得られること \\
\midrule
Perry and Wolf & ソフトウェアアーキテクチャ研究の基礎的な考え方を学べる。複数ビュー、アーキテクチャ様式、設計との違いなどの出発点を理解できる。 \\
Bass, Clements and Kazman & 品質特性、アーキテクチャ設計、記述、評価、技術的負債を実務的に学べる。 \\
Kruchten & 4+1 ビューモデルを学べる。論理、開発、プロセス、物理の四つのビューを、シナリオで結びつける考え方である。 \\
Rozanski and Woods & 利害関係者、関心事、ビュー、ビューポイントを使ってアーキテクチャを記述する実践的な方法を学べる。 \\
Taylor, Medvidović and Dashofy & アーキテクチャの基礎、理論、実践を包括的に学べる。実装、配置、非機能特性、信頼とセキュリティも扱う。 \\
Clements et al. & ビューとその種類を使って、アーキテクチャをどう文書化するかを学べる。 \\
Brown & 開発者の視点から、アーキテクチャドライバ、品質、制約、リスク管理を学べる。 \\
Fairbanks & リスク駆動で「必要十分な」アーキテクチャを行う考え方を学べる。軽量でアジャイルな文脈にも合う。 \\
Erder, Pureur and Woods & アジャイル、クラウド、DevOps 時代の継続的アーキテクチャを学べる。セキュリティ、回復性、拡張性にも重点がある。 \\
\bottomrule
\end{longtable}
